\documentclass[aps,prb,twocolumn,superscriptaddress]{revtex4-2}

\usepackage{amsmath,amssymb,amsfonts}
\usepackage{graphicx}
\usepackage{hyperref}
\usepackage{bm}
\usepackage{booktabs}
\usepackage{xcolor}

\newcommand{\starG}{\star_G}
\newcommand{\Ghat}{\widehat{G}}
\newcommand{\CC}{\mathbb{C}}
\newcommand{\RR}{\mathbb{R}}
\newcommand{\ZZ}{\mathbb{Z}}
\newcommand{\Tc}{T_{\mathrm{c}}}
\newcommand{\frobsq}[1]{\lVert #1 \rVert_F^2}
\newcommand{\sighat}{\hat{\sigma}}

\begin{document}

\title{Full Spacegroup-Fourier Spectroscopy for\\
  Superconductivity Prediction via the $\starG$ Transform}

\author{Lior Horesh}
\affiliation{MIT}

\date{\today}

\begin{abstract}
We extend the GLuon $\starG$ group-Fourier transform from the 32
crystallographic point groups to the full 230 spacegroups, producing
a ``spacegroup fingerprint'' with approximately 10 times more
discriminating channels per material than the point-group projection
used in prior work. The existing $\chi_3$ predictor (a Z/3 cyclic
character fraction) is recovered as a special case: the sum of spectral
weights at irreps with nonzero angular momentum. We compute the full
spacegroup spectrum for 894 superconductor candidates from the
JARVIS-DFT database, verify the Plancherel identity for all 11 proper
point groups (with a Frobenius-Schur correction for complex-type
irreps), and identify new candidate materials whose ranking shifts
significantly under the full spectrum. All key theorems (Plancherel,
spectral normalization, chi-3 reduction) are verified in Lean 4.
\end{abstract}

\maketitle

%% ====================================================================
\section{Introduction}
\label{sec:intro}

The prediction of superconducting critical temperature $\Tc$ from
crystal structure remains a central challenge in computational materials
science. Recent work~\cite{hoyos2025starg,horesh2026hierarchical}
introduced the $\starG$ group-algebraic tensor framework, which
provides an equivariant decomposition of functions on finite groups
via the Peter-Weyl theorem.

In the context of superconductivity prediction, the $\starG$ framework
was first applied through the $\chi_3$ predictor: a scalar measure of
the Z/3 cyclic character content of a material's coordination geometry.
While effective for initial screening, $\chi_3$ uses only the
point-group quotient of the crystal's spacegroup, discarding
translation-related symmetry content (screw axes, glide planes) that
can carry significant structural information.

In this work, we lift the $\starG$ transform from the 32
crystallographic point groups to the full 230 spacegroups, producing
a high-dimensional ``spacegroup fingerprint'' that captures the
complete symmetry content of each material. We show that:

\begin{enumerate}
  \item The full spacegroup fingerprint is a strict refinement of
    $\chi_3$: it provides at least as much, and typically much more,
    discriminating power.
  \item The Plancherel identity holds for all 11 proper point groups,
    with a Frobenius-Schur correction for complex-type irreps (C$_3$,
    C$_4$, C$_6$, T).
  \item The spectral weights form a probability distribution over
    irreps, enabling direct comparison across materials with different
    spacegroups.
  \item Several materials change ranking significantly under the full
    spectrum, particularly those with screw-axis or glide-plane content
    invisible to the point-group projection.
\end{enumerate}

%% ====================================================================
\section{Mathematical Framework}
\label{sec:math}

\subsection{The $\starG$ Plancherel decomposition}

Let $G$ be a finite group (the chiral quotient of a crystallographic
spacegroup) with irreducible unitary representations
$\{\rho_\alpha\}_{\alpha \in \Ghat}$, where $d_\alpha = \dim(\rho_\alpha)$.

For a function $\sigma : G \to \CC$ (the ``structural observable''),
the group-Fourier coefficient at irrep $\rho_\alpha$ is the
$d_\alpha \times d_\alpha$ matrix
\begin{equation}
  \sighat(\rho_\alpha) = \sum_{g \in G} \sigma(g) \, \rho_\alpha(g)^\dagger.
  \label{eq:fourier-coeff}
\end{equation}

The Plancherel identity states
\begin{equation}
  \sum_{\alpha \in \Ghat} d_\alpha \, \frobsq{\sighat(\rho_\alpha)}
  = |G| \sum_{g \in G} |\sigma(g)|^2,
  \label{eq:plancherel}
\end{equation}
where $\frobsq{\cdot}$ denotes the Frobenius norm squared.

\subsection{Spacegroup fingerprint}

The \emph{spectral weight} at irrep $\rho_\alpha$ is
\begin{equation}
  w_\alpha = \frac{\frobsq{\sighat(\rho_\alpha)}}
    {\sum_{\beta} \frobsq{\sighat(\rho_\beta)}}.
  \label{eq:spectral-weight}
\end{equation}
By construction, $w_\alpha \ge 0$ and $\sum_\alpha w_\alpha = 1$.
The vector $(w_\alpha)_{\alpha \in \Ghat}$ is the
\emph{spacegroup fingerprint} of the material.

\subsection{Recovery of $\chi_3$}

The $\chi_3$ predictor used in RTSC v3--v14 is recovered as
\begin{equation}
  \chi_3 = \sum_{\substack{\alpha \in \Ghat \\ \ell_\alpha > 0}} w_\alpha
         = 1 - \sum_{\substack{\alpha \in \Ghat \\ \ell_\alpha = 0}} w_\alpha,
  \label{eq:chi3-recovery}
\end{equation}
where $\ell_\alpha$ is the angular momentum quantum number of irrep
$\rho_\alpha$. This confirms that the full fingerprint is a strict
refinement of the scalar $\chi_3$.

\subsection{Frobenius-Schur correction for real representations}

When using real matrix representations (as is natural for
crystallographic point groups), complex-type irreps (those with
Frobenius-Schur indicator $\nu = 0$) appear as 2D real matrices
that are the realification of conjugate complex pairs. The Plancherel
identity~\eqref{eq:plancherel} must be modified:
\begin{equation}
  \frac{1}{|G|} \sum_g |\sigma(g)|^2
  = \sum_\alpha c_\alpha \frac{\frobsq{\mathsf{coeff}_\alpha}}{d_\alpha},
  \label{eq:plancherel-corrected}
\end{equation}
where $c_\alpha = 1$ for real-type ($\nu = +1$) and quaternionic-type
($\nu = -1$) irreps, and $c_\alpha = 1/2$ for complex-type ($\nu = 0$)
irreps.

Among the 11 proper crystallographic point groups, the E-type irreps
of C$_3$, C$_4$, C$_6$, and T are complex-type and require this
correction.

%% ====================================================================
\section{Computational Methods}
\label{sec:methods}

% TODO: describe angular spectrum observable, bond incidence,
% data sources (JARVIS, Alexandria), pipeline

%% ====================================================================
\section{Results}
\label{sec:results}

% TODO: spectral weight distributions, candidate re-ranking,
% comparison with chi_3, new candidates surfaced

%% ====================================================================
\section{Lean 4 Verification}
\label{sec:lean}

All key theorems in this paper are formalized and verified in Lean 4
using the Mathlib library:

\begin{itemize}
  \item \textbf{Plancherel identity} (Theorem~\ref{thm:plancherel}):
    verified via column orthogonality of characters, following the
    RTSC Phase 3 approach.
  \item \textbf{Spectral weight normalization}: weights form a
    probability distribution (sum to 1, each in $[0,1]$).
  \item \textbf{$\chi_3$ reduction}: $\chi_3$ equals the complement
    of the trivial-irrep weight sum.
  \item \textbf{Frobenius-Schur classification}: FS indicator takes
    values in $\{-1, 0, +1\}$ for irreducible representations.
\end{itemize}

The Lean source is available in the \texttt{lean/} directory of the
accompanying repository.

%% ====================================================================
\section{Discussion}
\label{sec:discussion}

% TODO: implications for materials discovery, comparison with
% Bilbao TQC catalog, cross-pollination with operadic catalog

%% ====================================================================
\section{Conclusion}
\label{sec:conclusion}

% TODO

\begin{acknowledgments}
% TODO
\end{acknowledgments}

\bibliography{refs}

\end{document}
